\documentclass[11pt,a4paper]{article}
\usepackage{hyperref}
\usepackage{natbib}
\usepackage{fontspec}
\usepackage{times}
\usepackage{xeCJK}

\setCJKmainfont{FandolSong}
\setCJKsansfont{AR PL UMing HK}
\setCJKmonofont{AR PL UKai HK}

\hypersetup{
    colorlinks=true,
    linkcolor=blue,
    urlcolor=blue,
}

\title{Writing Cantonese with XeLaTeX on Overleaf}
\author{LO, Chi-kiu (Jackie) {\tt 羅致翹}\\ Multilingual Text Processing \\ Digital Technologies Research Center \\National Research Council Canada (NRC-CNRC) \\ 1200 Montreal Road, Ottawa, ON K1A 0R6, Canada\\ {\tt chikiu.lo@nrc-cnrc.gc.ca } }
\date{June 2020}

\begin{document}
\maketitle
This article demonstrates how to write Cantonese characters in documents using XeLaTeX on Overleaf.

\section{Introduction}
The Hong Kong Supplementary Character Set (latest version HKSCS-2016 \citep{HKSCS2016}) is a set of specific Cantonese characters that is not included in common traditional Chinese character set for various reasons \citep{LuCheungXiongYuXu2014}. Without appropriate settings in the LaTeX document, these Cantonese characters may not be rendered properly. In this article, we show a working example on writing these characters with XeLaTex on Overleaf. 

\section{XeLaTeX and xeCJK}
To setup an Overleaf project for Chinese document using xeCJK with XeLaTeX, we follow the example in \url{https://www.overleaf.com/learn/latex/Chinese}. Essentially, it includes the following two steps:
\begin{enumerate}
\item Click on ``Menu'' at the top left corner and choose ``XeLaTeX'' as compiler in settings.
\item Add ``{\tt \textbackslash usepackage\{xeCJK\}}'' at the beginning of the document source.
\end{enumerate}

\section{Fonts}
The default xeCJK font on Overleaf is \emph{FandolSong}, which covers both Simplified Chinese and Traditional Chinese glyphs but not the characters in HKSCS-2016. For example, using \emph{FandolSong} as the CJK font, the second, fourth and the fifth character of 【石硤尾邨嘅居民】 cannot be rendered.

\citet{Lim2020} listed all the OTF/TTF fonts installed on Overleaf that can be used with fontspec in XeLaTeX. A general principle for selecting font when writing Cantonese is to use fonts specified with "HK" locality, such as \emph{AR PL UMing HK} and \emph{AR PL UKai HK}. As we can see below, using the two mentioned fonts, all the characters in previous example can be rendered properly. 

\begin{itemize}
    \item[] \emph{AR PL UMing HK}: {\sf【石硤尾邨嘅居民】} 
    \item[] \emph{AR PL UKai HK}:{\tt【石硤尾邨嘅居民】} 
\end{itemize}

To use the selected font, add ``{\tt \textbackslash setCJKmainfont\{<fontname>\}}'' at the beginning of the document source.

\bibliography{reference}
\bibliographystyle{acl_natbib}
\end{document}
